\section{Cap-free control of positive boundary layers}
\label{sec:v55-cap-free}
The finite-count error in a local thin-layer estimate is not a mass
which must be carried by the physical source. At small collision counts
there is a second estimate: the exact global mass of the same layer.
Combining these two inequalities removes the error uniformly in both
variables. This gives an explicit endpoint for higher integrability,
without using an exponential cap on the complete density.

We keep the original normalization
\[
 P(u)=m^2p_{n,R}(k,m,u),\qquad
 F_\varepsilon=m^2f^{\varepsilon,1}_{n,k,m,R},\qquad
 R_\varepsilon=P-F_\varepsilon\ge0.
\]
Here $m\ge2$, $1\le n\le m$, and $k\in\Z^2$. These are the
exact labels, not an occupation allowance. The source is
$\nu_R^*=c^{-1}\nu|_{Y_R^*}$. All roof inequalities are understood
almost everywhere and all local estimates are uniform in translated
intervals $I$ of fixed length $h>0$.

\subsection{Two finite estimates for the same source}
\begin{lemma}[Global all-margin mass]
\label{lem:v55-global-layer}
For $0<\varepsilon<\varepsilon_0<1$,
\begin{equation}\label{eq:v55-global-layer}
 \sup_R\sum_{n=1}^m\sum_{k\in\Z^2}
        \int_{\R} R_{\varepsilon,n,k,m,R}(u)\,du
                         \le C m^2\varepsilon.
\end{equation}
The same bound holds for either positive physical component, and with
an extra factor $M$ in variation for an insertion of modulus at most $M$.
\end{lemma}
\begin{proof}
Every factor of the guard $\Xi^\varepsilon$ is between zero and one.
Its complement is bounded by the sum of the strip indicators used in
Lemma~\ref{lem:v54-fixed-band-windows}. At collision $j$, an incidence
strip of width $a$ has collision mass at most $Ca^2$, and a section
strip of width $a$ has mass at most $Ca$. The latter follows directly
from the bounded perimeter of the section rectangles in the flat
collision coordinates. A clearance strip has mass at most $Ca$;
it is read at collision $j+1$ as in
Section~\ref{sec:v49-physical-layers}. These are invariant one-flight
mass estimates, not conditional independence assertions.

Use $a=2\varepsilon q_0^{d_j/4}$ for contacts and section decisions,
and $a=2\varepsilon q_0^{d'_j/4}$ for flights, where
$q_0=47/53$, $d_j=\min(j,m-j)$ and
$d'_j=\min(j,m-1-j)$. At most two indices have any given depth.
The convergent geometric sums give
$\int(1-\Xi^\varepsilon)\,d\nu\le C(\varepsilon+\varepsilon^2)$.
Enlarging the section source costs $1/c$ once. For fixed $m$, the
original events $\Pi_{n,k,m,R}$ are disjoint as $(n,k)$ varies.
Consequently summing their pushforward masses and multiplying by $m^2$
proves \eqref{eq:v55-global-layer}. Neither the endpoint membership
at $m$ nor this summation adds a term to the occupation at
$0,\ldots,m-1$. The positive physical components are dominated by
$R_\varepsilon$, since $\Xi^\varepsilon\le H^\varepsilon$.
Source domination proves the weighted assertion.
\end{proof}

\begin{lemma}[Removal of a finite-count error]
\label{lem:v55-two-regimes}
Let $0<\alpha<\beta$, $\kappa>0$ and $r\ge0$. Suppose that
$U_m(\varepsilon)\ge0$ satisfies, simultaneously for all $m\ge2$
and $0<\varepsilon<\varepsilon_0<1$,
\[
 U_m(\varepsilon)\le C_1(\varepsilon^\alpha+e^{-\kappa m}),
 \qquad U_m(\varepsilon)\le C_2 m^r\varepsilon^\beta.
\]
Then $U_m(\varepsilon)\le C\varepsilon^\alpha$ with a constant
independent of both variables.
\end{lemma}
\begin{proof}
If $\kappa m\ge\alpha\log(1/\varepsilon)$, use the first inequality.
Otherwise the second inequality gives
\[
 U_m(\varepsilon)\le
 C_2(\alpha/\kappa)^r[\log(1/\varepsilon)]^r\varepsilon^\beta
 \le C\varepsilon^\alpha,
\]
because $x^r e^{-(\beta-\alpha)x}$ is bounded for $x\ge0$.
For $r=0$ the same calculation has the factor $x^0=1$.
Both estimates concern the same source at the same finite count;
no asymptotic diagonal or exchange of limits is used.
\end{proof}

\begin{theorem}[Simultaneous local thin mass without a spectral error]
\label{thm:v55-error-free-mass}
For every interval $I$ of length $h>0$, all $m\ge2$ and
$0<\varepsilon<\varepsilon_0$,
\begin{equation}\label{eq:v55-error-free-mass}
 \sup_{R,n,k}\int_I R_\varepsilon(u)\,du
                       \le C(1+h)\varepsilon^{1/16}.
\end{equation}
The same bound holds for each normalized positive incidence and
clearance density and for their sum. For an arbitrary bounded
insertion, its variation is bounded by $M$ times this right side.
The protection scale is permitted to depend on $m$.
\end{theorem}
\begin{proof}
In \eqref{eq:v54-all-margin-window} absorb
$m^3\rho^{m/2}$ into $Ce^{-\kappa m}$ with
$\kappa=-\frac14\log\rho>0$. Divide the local integral by $1+h$.
Its second bound is \eqref{eq:v55-global-layer}, since a local mass
is no larger than the global mass and $1+h\ge1$. Apply
Lemma~\ref{lem:v55-two-regimes} with
$(\alpha,\beta,r)=(1/16,1,2)$. Domination gives the other assertions.
The fixed auxiliary band used to establish the first inequality has
not been altered. The dependence of the final constant on its decay
rate is allowed; the exponent $1/16$ does not have that dependence.
\end{proof}

\subsection{A weak endpoint independent of a height cap}
We use the weak $L^p$ distribution convention
$\|f\|_{p,\infty;I}^p=\sup_{L>0}L^p
 |\{u\in I:|f(u)|>L\}|$. This specifies a quasi-norm, not a
claim about the strong $L^p$ integral.

\begin{theorem}[A positive-source endpoint principle]
\label{thm:v55-endpoint-principle}
Suppose $P_m\ge0$ has exact positive decompositions
$P_m=A_{m,\varepsilon}+B_{m,\varepsilon}$ on a finite measure
space, and, uniformly for all finite $m,\varepsilon$,
\[
 \|A_{m,\varepsilon}\|_\infty\le C_A\varepsilon^{-K},\quad
 \int B_{m,\varepsilon}\le C_1(\varepsilon^\alpha+e^{-\kappa m}),
 \quad \int B_{m,\varepsilon}\le C_2m^r\varepsilon^\beta,
\]
where $K,\kappa>0$ and $0<\alpha<\beta$. Suppose also that
$\int P_m$ is uniformly bounded. Put $a=\alpha/K$ and $p_c=1+a$.
Then, for a fixed $L_0$ and all $L\ge L_0$,
\begin{equation}\label{eq:v55-abstract-tail}
 \int_{\{P_m>L\}}P_m\le CL^{-a},\qquad
 |\{P_m>L\}|\le CL^{-p_c}.
\end{equation}
Thus $P_m$ is uniformly in weak $L^{p_c}$ and in strong $L^q$ for
every $1<q<p_c$. No complete-source height bound is assumed.
\end{theorem}
\begin{proof}
Lemma~\ref{lem:v55-two-regimes} first gives
$\int B_{m,\varepsilon}\le C\varepsilon^\alpha$. Choose
$\varepsilon=(2C_A/L)^{1/K}$, with $L_0$ large enough that this is
admissible. On $\{P_m>L\}$ one has
$B_{m,\varepsilon}\ge P_m/2$. This proves the first bound in
\eqref{eq:v55-abstract-tail}; division by $L$ proves the second.
For $1<q<1+a$, the layer-cake integral
\[
 \int P_m^q\le\int P_m+(q-1)\int_1^\infty
 L^{q-2}\int_{\{P_m>L\}}P_m\,dL
\]
uses $P_m^q\le P_m$ on $P_m<1$.
The tail is integrable because $q-1<a$. For smaller thresholds the
uniform first moment bounds the weak quasi-norm. All integration now
uses an error-free tail, so there is no exponential cap or rate ratio
in the exponent range.
\end{proof}

For the actual record define the explicit endpoint
\begin{equation}\label{eq:v55-explicit-endpoint}
 a_c=\frac1{144},\qquad p_c=\frac{145}{144}.
\end{equation}
\begin{corollary}[The complete exact-label density at the weak endpoint]
\label{cor:v55-physical-endpoint}
For all translated intervals $I$ of length $h>0$,
\begin{align}
 \sup_{R,n,k,m}\int_{I\cap\{P>L\}}P(u)\,du
       &\le C(1+h)L^{-1/144} &&(L\ge L_0),
                                      \label{eq:v55-physical-tail}\\
 \sup_{R,n,k,m}\|P\|_{p_c,\infty;I}^{p_c}
       &\le C(1+h),\qquad
 \sup_{R,n,k,m}\int_I P^q\,du\le C_q(1+h)
                  &&(1<q<p_c).       \label{eq:v55-physical-weak}
\end{align}
The supremum includes every $m\ge2$ and $1\le n\le m$.
\end{corollary}
\begin{proof}
Apply Theorem~\ref{thm:v55-endpoint-principle} with the original guard
split, $K=9$, $\alpha=1/16$, $\beta=1$ and $r=2$.
Proposition~\ref{prop:v54-protected-height} supplies the finite
protected height. Lemma~\ref{lem:v54-fixed-band-windows} and
Lemma~\ref{lem:v55-global-layer} supply the two mass bounds and
uniform first moments. Equivalently use the already error-free
Theorem~\ref{thm:v55-error-free-mass}. Divide the measure on $I$ by
$1+h$ before applying the principle, and multiply back afterwards.
The complete finite-count density exists by
Theorem~\ref{thm:v44-complete-density-bound}, but its exponential
height estimate is not an input to this argument.
\end{proof}

\paragraph{The remaining pointwise estimate.}
The endpoint is independent of $\kappa/\Gamma$, unlike the sufficient
interior exponent in \eqref{eq:v54-admissible-exponent}. Constants
still depend on the fixed-band estimates. Weak $L^{145/144}$ does
not assert strong integrability at that endpoint, and neither assertion
is an essential-height bound. In particular the positive incidence
and clearance criterion in Corollary~\ref{cor:v51-positive-criterion}
remains the original pointwise target. No physical seam has been
crossed in deriving the new estimates.
